% Section 7: exact sampling and speculative decisions (paper; theory from
% PR #7, the P1 and P3 analyses listed in TASKS.md, and the earlier manuscript).
% Claims and status:
%  - Theorems race/fixednoise/guards/residualrace/residualself, Propositions rng/coupling/
%    commonmass/boundedrange/kl/frontier: proved here (finding). Known principles are
%    attributed in the text.
%  - Synthetic drift work counts: finding on main (data/synthetic_drift.csv,
%    evidence/synthetic_drift.json), constructed matrices, not timings.
%  - Race candidates at T > 0 and fixed-noise acceptance on real drafts: pending.
\section{Exact sampling and speculative decisions}\label{sec:decisions}
Greedy selection is the easy consumer. This section turns the same evidence into sampled decisions: a token drawn from the target law, an accept or reject bit for a drafted token, and the residual or bonus sample that follows. The evidence may come from transport or from self-evidence; every statement below takes enclosures of the logits or of their exponentials as input and does not care which mechanism produced them.

\subsection{A certified race under a keyed noise field}\label{sec:race}
\begin{theorem}[Certified race]\label{thm:race}
For any fixed noise vector $g$ and $T>0$, the ladder of Theorem~\ref{thm:ladder} with $f_i(x)=x/T+g_i$ returns the S-race token, which is the token of a dense race with the same $g$ and tie rule. If the $g_i$ are i.i.d.\ standard Gumbel and independent of $z$, the token has law $\softmax(z/T)$.
\end{theorem}
\begin{proof}
Theorem~\ref{thm:ladder} and the Gumbel-max identity~\cite{gumbel1954,papandreou2011}; ties have probability zero under continuous noise.
\end{proof}
Using bounds to avoid evaluating most perturbed scores is the principle of A* Sampling~\cite{astar} and of lazy Gumbel sampling from a retrieved top set~\cite{mussmann2017}. What the fixed field adds is pathwise equality with the dense race, which a lazily instantiated field does not give. Top-$k$ sampling applies the ladder to membership (drop $i$ when $hi_i$ is below the $k$-th largest $lo$) and races over the certified set; equality then also needs the engine's tie semantics at the $k$-th position. Min-$p$ admits $i$ when $z_i\ge z_{\max}+T\log(\text{min-}p)$, which intervals decide without a partition function. Top-$p$ admission depends on the mass above each token relative to $Z$, needs partition bounds, and is usually cheaper to hand to the ordinary path.

\begin{proposition}[RNG contract]\label{prop:rng}
(i) Seeded equality between a certified and a dense procedure holds when both evaluate the same function of the same noise values with the same tie rule. This requires a counter-based field, a pure function of (seed, request, absolute position, draw purpose, token) that the certified path can evaluate for any token in any order; a stream consumed in evaluation order does not qualify, and in the exact reference assigning a stream to rows in candidate order changed the winner in 247 of 300 cases (\nolinkurl{evidence/precision/precision_tests.json}). (ii) Equality in distribution needs i.i.d.\ noise with the correct marginal, independent of the proposal draw, the acceptance uniform and the scores; a residual race also needs independence from the acceptance event. (iii) Fixed-noise verification (Section~\ref{sec:fixednoise}) needs the noise of a position to be the same whether the position is drafted or decoded plainly, so it must be keyed by absolute position, not by cycle or batch slot. (iv) A hash with $2^{32}$ levels makes the law exact only for the idealized continuous field; the stock seeded sampler shares that approximation.
\end{proposition}
The stock seeded sampler (Section~\ref{sec:stock}) already has the form (i) requires. FlashSampling's noise offset depends on the tiling~\cite{flashsampling}, so it cannot serve as a certified path's field without change.

\subsection{How much of a sample can be emitted before the scores are known}\label{sec:commonmass}
A race needs the winner's interval to separate. A different construction can emit a token with some probability before any candidate is separated, and then finish the rest of the law after more computation. Let $\mathcal P(I)$ be the set of target laws compatible with the information $I$ held so far.

\begin{proposition}[Common mass]\label{prop:commonmass}
Consider procedures that, before acquiring more information, emit token $i$ with probability $r_i$ (the same for every $p\in\mathcal P(I)$) and otherwise continue until $p$ is known. Such a procedure outputs exactly $p$ for every $p\in\mathcal P(I)$ if and only if $r_i\le\inf_{p\in\mathcal P(I)}p_i$ for every $i$, and then it may finish by sampling the residual $(p-r)/(1-\alpha)$ with $\alpha=\sum_ir_i$. The largest immediate probability is
\[\alpha^\ast(I)=\sum_i\inf_{p\in\mathcal P(I)}p_i.\]
\end{proposition}
\begin{proof}
If the procedure is exact, the final probability of $i$ is at least the probability $r_i$ of emitting it immediately, for every compatible $p$; so $r_i\le\inf_pp_i$. Conversely, if $r_i\le\inf_pp_i$ then $p-r\ge0$ and sums to $1-\alpha$ for every compatible $p$, so the residual is a law, and the output law is $r+(1-\alpha)\cdot(p-r)/(1-\alpha)=p$.
\end{proof}
For two candidate laws $\mathcal P=\{p,q\}$, $\alpha^\ast=\sum_i\min(p_i,q_i)=1-\TV(p,q)$, the overlap of a maximal coupling, which is exactly the acceptance probability of speculative sampling~\cite{specdecode,chen2023specsampling}. The proposition says the same idea applies to any set of laws the evidence leaves open, and that nothing can do better without more information.

A two-token example shows that exact sampling can finish where greedy selection cannot. Suppose the evidence only says that the probability of token~1 lies in $[0.49,0.51]$. Then $\alpha^\ast=0.49+0.49=0.98$. Draw $U$ uniform on $[0,1)$: emit token~1 if $U<0.49$, token~0 if $U\ge0.51$, and otherwise compute $p$ and emit token~1 exactly when $U<p$, reusing the same $U$. The output has law $p$, and 98\% of draws never need the scores; the greedy decision, by contrast, is undecided because the interval contains $1/2$. For two tokens this procedure is inverse-CDF sampling with bounds on the CDF. In general it gives equality in distribution, not seeded equality with a dense race.

\begin{proposition}[Bounded logit error]\label{prop:boundedrange}
Let $p=\softmax(z/T)$ and $q=\softmax(\tilde z/T)$, and suppose $\max_i(z_i-\tilde z_i)-\min_i(z_i-\tilde z_i)\le\Delta$. Then $p_i\ge e^{-\Delta/T}q_i$ for every $i$. Consequently, with $\alpha=e^{-\Delta/T}$:
\begin{enumerate}
\item sampling $q$ with probability $\alpha$, and otherwise computing $p$ and sampling $(p-\alpha q)/(1-\alpha)$, samples $p$ exactly;
\item always sampling $q$ gives $\TV(p,q)\le1-\alpha$.
\end{enumerate}
\end{proposition}
\begin{proof}
Write $d_i=z_i-\tilde z_i$ with $m\le d_i\le M$ and $M-m\le\Delta$. Then $p_i/q_i=e^{d_i/T}/\sum_jq_je^{d_j/T}\ge e^{m/T}/e^{M/T}\ge e^{-\Delta/T}$. Part~1 is Proposition~\ref{prop:commonmass} with $r=\alpha q$. For part~2, $\TV(p,q)=\sum_i(q_i-p_i)_+\le\sum_i(1-\alpha)q_i=1-\alpha$.
\end{proof}
With self-evidence the range is at most twice the largest envelope half-width, so $\alpha\ge e^{-2\max_i\beta_i/T}$. The trade-off against the race is different: this sampler needs the cheap pass's full softmax and, with probability $1-\alpha$, the whole exact head, whereas the race re-scores a candidate set whose expected size grows like $e^{4\beta/T}$ and returns the seeded token.

\begin{proposition}[Approximate mode: KL and sequence total variation]\label{prop:kl}
Suppose $\min_c\norm{z-\tilde z-c\mathbf 1}_\infty\le\varepsilon$ at temperature one (at temperature $T$ replace $\varepsilon$ by $\varepsilon/T$). Then $\KL(p\,\|\,q)\le\varepsilon^2/2$ and $\KL(q\,\|\,p)\le\varepsilon^2/2$, so $\TV(p,q)\le\varepsilon/2$. If an approximate decoder samples $q_t$ in place of $p_t$ at every position and the bound holds with $\varepsilon_t$ at every reachable prefix, the laws $P$ and $Q$ of the first $N$ tokens satisfy $\KL(P\,\|\,Q)\le\tfrac12\sum_t\varepsilon_t^2$ and
\[\TV(P,Q)\le\tfrac12\Bigl(\sum_{t=1}^N\varepsilon_t^2\Bigr)^{1/2}.\]
\end{proposition}
\begin{proof}
Let $d=z-\tilde z-c\mathbf 1$ with $\norm d_\infty\le\varepsilon$; the shift $c$ changes neither law. Then $p$ is $q$ tilted by $d$: $p_i=q_ie^{d_i}/\E_q[e^d]$. Let $\psi(\lambda)=\log\E_q[e^{\lambda d}]$, so $\psi(0)=0$ and $\KL(p\,\|\,q)=\E_p[d]-\psi(1)=\psi'(1)-\psi(1)=\int_0^1\lambda\psi''(\lambda)\,d\lambda$. Here $\psi''(\lambda)$ is the variance of $d$ under the tilted law $q^{(\lambda)}$, at most $(2\varepsilon)^2/4=\varepsilon^2$ because $d$ takes values in an interval of length $2\varepsilon$. Hence $\KL(p\,\|\,q)\le\varepsilon^2/2$; the reverse divergence follows by exchanging the roles of $p$ and $q$. Pinsker's inequality gives $\TV\le\sqrt{\KL/2}\le\varepsilon/2$. For sequences, the chain rule gives $\KL(P\,\|\,Q)=\sum_t\E_P[\KL(p_t\,\|\,q_t)]\le\sum_t\varepsilon_t^2/2$, and Pinsker's inequality again gives the total-variation bound.
\end{proof}
To guarantee a sequence total variation of at most $0.01$ over $N=1{,}024$ tokens with a uniform bound, the proposition needs a logit error, up to a common shift, of at most $\varepsilon=6.25\times10^{-4}$, far tighter than any int8 envelope in Section~\ref{sec:selfevidence}. An approximate head is therefore a quality trade-off to be measured (Section~\ref{sec:stack}); these bounds cannot certify it as harmless over long outputs.

\subsection{Fixed-noise verification}\label{sec:fixednoise}
\begin{theorem}[Fixed-noise verification]\label{thm:fixednoise}
Let $G_{t,i}$ be keyed by (seed, absolute position $t$, token $i$) and i.i.d.\ standard Gumbel across $(t,i)$. Define $Y_t=\min\argmax_i\bigl(z_{t,i}(Y_{<t})/T+G_{t,i}\bigr)$. A speculative decoder that accepts a drafted token at position $t$ if and only if it equals $\argmax_i(z_{t,i}(x_{<t})/T+G_{t,i})$ at the drafted prefix, and emits that argmax at the first mismatch, outputs $Y$ for every drafter and every depth policy, and $Y$ has the target's autoregressive law.
\end{theorem}
\begin{proof}
By induction on positions, every accepted token equals the argmax at a prefix equal to $Y_{<t}$, so it is $Y_t$; at the first mismatch the prefix is $Y_{<t}$ and the emitted token is $Y_t$. The drafter only changes how many positions one target pass verifies. Given $Y_{<t}$, $G_{t,\cdot}$ is independent of $z_t(Y_{<t})$, so $Y_t\mid Y_{<t}\sim\softmax(z_t(Y_{<t})/T)$.
\end{proof}
This is the drafter-invariant construction of Daliri, Musco and Suresh~\cite{daliri2025}. For the certified head it matters because verification becomes a perturbed argmax, which Theorem~\ref{thm:race} certifies with no partition function and no residual sampler. The theorem treats $z_t$ as a function of the prefix alone. In the engine the verify pass computes the logits at a different batch shape from plain decoding, and Section~\ref{sec:divergence} shows that stock configurations at different shapes can disagree, so what a fixed-noise verifier reproduces token for token is the seeded sampler applied to the verify pass's own logits; equality with seeded plain decoding needs the verify-shape and plain-decode logits to agree as well, for example under batch-invariant kernels.

\begin{proposition}[Acceptance under fixed-noise coupling]\label{prop:coupling}
If the drafter proposes $\argmax_i(\log q_i+G_i)$ and the target accepts if and only if it equals $\argmax_i(\log p_i+G_i)$, then
\[P(\text{accept})=\sum_i\Bigl(\sum_j\max\bigl(\tfrac{p_j}{p_i},\tfrac{q_j}{q_i}\bigr)\Bigr)^{-1},\qquad
\tfrac12\sum_i\min(p_i,q_i)\le\sum_i\frac{p_iq_i}{p_i+q_i}\le P(\text{accept})\le\sum_i\min(p_i,q_i).\]
\end{proposition}
\begin{proof}
Both choose $i$ if and only if $G_i>\max_{j\ne i}(G_j+m_j)$ with $e^{m_j}=\max(p_j/p_i,q_j/q_i)$, an event of probability $(1+\sum_{j\ne i}e^{m_j})^{-1}$. The bounds follow from $\max(x,y)\le x+y$, from $\max(x,y)\ge x$ and $\ge y$ summed over $j$, and from $p+q\le2\max(p,q)$.
\end{proof}
Per position, fixed-noise acceptance is at most the $\min(1,p/q)$ rate $1-\TV(p,q)$ and at least half of it; the two agree when $p=q$. Whether the cheaper certificate pays for any lost acceptance is an empirical question \pending{acceptance under fixed-noise coupling against standard verification on real MTP and DFlash drafts at $T=1$ and $T=0.7$}.

\paragraph{The sampling contract adopted here.} The certified speculative path uses fixed-noise verification with exactly the stock seeded sampler's noise: $G_{t,i}=-\log(-\log x_{t,i})$ in FP64, with $x_{t,i}$ the 32-bit MurmurHash of (seed, absolute position $t$, token $i$) scaled to $[0,1]$, as in \code{multinomial\_with\_seed} (Section~\ref{sec:stock}). A drafted token is accepted if and only if it equals the target's keyed sample at its position, and at the first mismatch that sample is emitted. The output is equal in law to target sampling, and it is drafter-invariant: the same keyed noise and the same verifier logits give the same tokens for any drafter and any depth policy. Three further consequences follow. For SGLang's default point-mass drafts the per-position acceptance probability is unchanged: the stock rule accepts a draft $d$ when a uniform coin falls below $p(d)$, and the keyed rule accepts it when $d$ wins the race, which also has probability $p(d)$. With a seed, every committed token is the stock seeded sampler's token for the verify pass's own stock logits at that position (Theorem~\ref{thm:fixednoise} under R-stock at the verify shape), provided that every row whose race the certificate cannot decide under the stock rounding of $\log\softmax$ runs the stock sampler. The output equals seeded decoding without speculation only where the verify-pass and plain-decode logits agree, which the engine does not guarantee because the two use different kernels and shapes; the divergence analysis of Section~\ref{sec:divergence} is the measurement that will say how often they differ; stock speculative sampling, by contrast, is equal to non-speculative decoding only in law, because its coins and its inverse-CDF residual consume randomness differently. Without a seed the output equals stock in law. Top-$p$ requests fall back to the stock path. The acceptance guards below are analysis and are not built: under the stock-token contract the stock inverse-CDF residual would still need a dense head row per verify step. The design is decided; its measurement is \pending{kernel: token-for-token agreement with the stock seeded sampler on the verify pass's logits for replayed steps, and acceptance against the stock speculative sampler}.

\subsection{Acceptance guards from incomplete mass}\label{sec:guards}
For one reached proposal $X=x$ write $w_x=e^{z_x}$ and $Z=\sum_ie^{z_i}$, and draw $U\in[0,1)$ with the reference convention. For $q_x>0$, acceptance is equivalent to
\begin{equation}Uq_xZ<w_x,\label{eq:accept}\end{equation}
and the $\min(1,\cdot)$ need not be evaluated because $U<1$ already handles $p_x\ge q_x$.

\begin{theorem}[Acceptance guards with enclosures]\label{thm:guards}\label{thm:guard}
Let $w_x\in[N^-,N^+]$ and $Z\in[Z^-,Z^+]$ with $N^->0$, let $0<q_x\le1$ be the actual proposal mass and $U\in[0,1)$. Then $Uq_xZ^+<N^-$ implies acceptance, $Uq_xZ^-\ge N^+$ implies rejection, the two never both hold, and replacing an enclosure by a tighter one cannot reverse a certified decision. For $U$ uniform and independent of the enclosures, neither guard fires with probability
\[P_?=\min\bigl(1,N^+/(q_xZ^-)\bigr)-\min\bigl(1,N^-/(q_xZ^+)\bigr).\]
If $N^-=N^+=w_x$ and every row outside a re-scored set $R$ has a mass enclosure whose upper and lower ends differ by a factor at most $\kappa$, then with $\pi=\sum_{i\notin R}p_i$,
\[P_?\le\frac{p_x}{q_x}\cdot\frac{(\kappa-1)\pi}{1-\pi(1-1/\kappa)}.\]
\end{theorem}
\begin{proof}
The implications are transitivity with~\eqref{eq:accept}; both firing would give $N^+\le Uq_xZ^-\le Uq_xZ^+<N^-$. A tighter enclosure stays on the same side of the true values. The unresolved uniforms form $[N^-/(q_xZ^+),N^+/(q_xZ^-))\cap[0,1)$. For the bound, $P_?\le(w_x/q_x)(Z^+-Z^-)/(Z^+Z^-)$; $Z^+-Z^-\le\sum_{i\notin R}e^{z_i}(\kappa-1)$ because each lower mass is at most $e^{z_i}$; and $Z^-\ge Z-\sum_{i\notin R}e^{z_i}(1-1/\kappa)$ because each lower mass is at least $e^{z_i}/\kappa$. Divide by $Z$.
\end{proof}
These guards are analysis in this revision: the engine path adopted above uses fixed-noise verification, which needs no partition function, and a certified version of the stock rule would still pay for a dense residual row. Transported tile masses (Theorem~\ref{thm:transport}) and exponentiated self-evidence intervals (Theorem~\ref{thm:envelope}) both supply $Z^\pm$. With self-evidence $\kappa=e^{2\hat\beta}$ times the outward rounding of the exponential; at $\beta=1$ ($\kappa\approx7.4$), keeping $P_?$ below 10\% at $p_x/q_x\approx1$ requires the unre-scored mass $\pi$ to be below about 1.5\%. After an adaptive refinement chosen using $U$, the probability statement no longer applies, although the guards remain sound pointwise. The Lean formalization proves the guards with enclosed numerators and their exclusivity (\code{formal/CertifiedArgmax.lean}).

\subsection{Residual and bonus work}\label{sec:residual}
A rejected proposal still requires a sample from $(p-q)_+$, and knowing that rejection occurred does not supply one. When all drafted tokens are accepted, a bonus token needs a target sample. A certificate that decides acceptance cheaply and then completes the whole head for the residual saves little; Section~\ref{sec:synthetic} shows this happening. Two facts about the engine and one construction reduce the problem.

First, SGLang's default sampled verification uses point-mass proposals (Section~\ref{sec:stock}). With $q$ a point mass at $x$, the residual is the target restricted to the untried tokens, a race over $V\setminus\{x\}$ that needs no partition function. Second, under fixed-noise verification there is no residual at all. Third, for a proposal with small support there is an exact bounded race.

Let $Q=\{i:q_i>0\}$. With unnormalized target weights $w_i$ and partition $Z$, the residual can be sampled from weights
\begin{equation}
r_i=(w_i-Zq_i)_+,\qquad\frac{r_i}{\sum_jr_j}=\frac{(p_i-q_i)_+}{\sum_j(p_j-q_j)_+},\label{eq:residualweights}
\end{equation}
and $r_i=w_i$ outside $Q$. Given $Z^-\le Z\le Z^+$, positive-part monotonicity gives $(w_i-Z^+q_i)_+\le r_i\le(w_i-Z^-q_i)_+$ inside $Q$.

\begin{theorem}[Bounded residual race]\label{thm:residualrace}
Assume valid enclosures of the target weights and of $Z$, the actual proposal law, a nonzero residual, and a priority field $\pi_i=e^{G_i}>0$ that is fixed and independent of the proposal draw, the acceptance uniform and the rejection event. A refinement procedure that returns only an interval-certified ordered maximizer of $r_i\pi_i$, or completes the dense race, returns the same token as the dense race. With i.i.d.\ standard Gumbel $G_i$ it samples the residual law~\eqref{eq:residualweights} exactly.
\end{theorem}
\begin{proof}
The enclosures of $r_i$ multiplied by $\pi_i>0$ enclose the priorities. A certified maximizer dominates every possible competitor, including the actual one, and a completed evaluation yields the dense ordered maximizer. The distributional statement is Gumbel-max applied to the nonzero residual weights.
\end{proof}

\begin{theorem}[Residual races with self-evidence]\label{thm:residualself}
Replace the enclosures above by $r_i\in[(N^-_i-Z^+q_i)_+,(N^+_i-Z^-q_i)_+]$ for $q_i>0$ and $r_i\in[N_i^-,N_i^+]$ for $q_i=0$. The conclusions of Theorem~\ref{thm:residualrace} hold. For a point-mass proposal at $x$, $r_x=0$ as soon as $Z^-\ge N^+_x$, and the residual is a race over $V\setminus\{x\}$ that needs no partition bound. SGLang's target-only rule (accept sibling $j$ if $UZ<\sum_{l\le j}w_{x_l}$; otherwise sample the target restricted to untried tokens) is certified by the guards with a cumulative numerator and an exclusion race.
\end{theorem}
\begin{proof}
Positive-part monotonicity gives the enclosures; the rest is the proof of Theorem~\ref{thm:residualrace}. For the point mass, $w_x-Zq_x=w_x-Z\le0$.
\end{proof}
The engine samples that residual by inverse CDF, so a race matches it only in distribution. For a support chosen after the draft pass, retaining the top $K+1$ perturbed entries per tile tightens the outside-support maximum after excluding up to $K$ support tokens; it is not needed for validity, since $r_i\le w_i$ makes the plain tile maximum an upper bound. For a bonus sample the construction is simpler still: a certified race over the target scores with the bonus position's field, with no partition function.

\subsection{A cheap decision is not a cheap verification}\label{sec:synthetic}
The inherited CPU experiment makes the last point concrete. It constructs two families of $128\times8$ integer heads with 16-row tiles and 36 paired fixtures each: one clusters rows around separated centres, the control uses unstructured rows. For each fixture $h^t=h^d+d\,v$ with drift scale $d\in\{0,1,2,4\}$; proposals are sampled from the anchor head's law, and all probability arithmetic is exact rational. The experiment counts distinct target rows projected, which is work, not time.

At zero drift the summaries describe the target exactly and the acceptance decision needs only its candidate row, an identity check rather than a forecast. At nonzero drift the unstructured head defeats greedy screening. More subtly, large drift makes rejection easy to certify, which lowers the rows needed for the decision while destroying acceptance itself, and completing the residual then erases the saving (Figure~\ref{fig:synthetic}). At drift four the clustered family decides with 11.4\% of the rows but needs 96.5\% once the dense residual is charged. The experiment does not refute transport on trained heads; it refutes the inference that a cheap acceptance decision means cheap verification. The residual races above are the answer it motivates, and they remain to be measured.

\begin{figure}[tbp]
\centering
\begin{tikzpicture}
\begin{axis}[width=.66\linewidth,height=5.4cm,xmin=-.2,xmax=4.2,ymin=0,ymax=1.05,xtick={0,1,2,4},
  xlabel={Drift scale $d$},ylabel={Fraction of target rows projected},
  grid=major,grid style={draw=ink!10},legend style={at={(1.02,.5)},anchor=west,draw=none,font=\scriptsize,cells={align=left}},legend cell align=left]
\addplot[thick,teal,mark=*,discard if not={family}{clustered},unbounded coords=discard] table[x=drift,y=accept_projected_fraction,col sep=comma]{../data/synthetic_drift.csv};
\addlegendentry{Clustered: decision}
\addplot[thick,teal,dashed,mark=o,discard if not={family}{clustered},unbounded coords=discard] table[x=drift,y=decision_plus_residual_projected_fraction,col sep=comma]{../data/synthetic_drift.csv};
\addlegendentry{Clustered: decision\\+ dense residual}
\addplot[thick,ochre,mark=square*,discard if not={family}{unstructured},unbounded coords=discard] table[x=drift,y=accept_projected_fraction,col sep=comma]{../data/synthetic_drift.csv};
\addlegendentry{Unstructured: decision}
\addplot[thick,ochre,dashed,mark=square,discard if not={family}{unstructured},unbounded coords=discard] table[x=drift,y=decision_plus_residual_projected_fraction,col sep=comma]{../data/synthetic_drift.csv};
\addlegendentry{Unstructured: decision\\+ dense residual}
\addplot[thin,muted,mark=triangle,discard if not={family}{clustered},unbounded coords=discard] table[x=drift,y=acceptance_fraction,col sep=comma]{../data/synthetic_drift.csv};
\addlegendentry{Clustered: accepted}
\end{axis}
\end{tikzpicture}
\caption{Work counts on constructed heads, not timings (\nolinkurl{data/synthetic_drift.csv}, \nolinkurl{evidence/synthetic_drift.json}). Solid lines count the rows needed to decide acceptance; dashed lines add a dense residual completion after each rejection. The thin line is the fraction of proposals accepted in the clustered family. Evidence production, bounds, bonus sampling and the backbone are excluded.}
\label{fig:synthetic}
\end{figure}

\subsection{A demand frontier over a fixed protocol}\label{sec:frontier}
For a chain of $L$ proposals each head row is accepted, rejected or unresolved. The first certified rejection at position $r$ proves that no row after $r$ can affect the cycle's output, while earlier unresolved rows still matter because one of them may be the true first rejection.

\begin{proposition}[Safe suppression after a certified rejection]\label{prop:frontier}
If every issued certificate equals the decision of a fixed reference verifier on its row, the first true rejection is no later than the first certified rejection. Suppressing head work strictly after the latter does not change the reference output, provided completion and state commit use the actual first rejection.
\end{proposition}
\begin{proof}
At the certified rejection the reference decision is rejection; an earlier unresolved row may also reject, and no later row can be the first rejection. All accepted outputs, the replacement token and the state prefix are determined by rows up to the first rejection.
\end{proof}
The statement is for chains. SGLang's trees share one coin across siblings and choose the next coin by the accepted node, so a tree frontier must follow the tree's paths. The target backbone has already run over every admitted position, so the frontier saves head work, not transformer layers.

\subsection{When adaptive verification changes the sampling law}\label{sec:stopping}
Avoiding work whose result is already determined differs from choosing which sampled proposals count. DSpark states that lossless speculation requires admission decisions that do not depend on future candidate tokens, and D-cut gives a counterexample and a repair~\cite{dspark,dcut}. A two-token illustration: let $p=q=(1/2,1/2)$, draw $X\sim q$, discard it and draw afresh from $p$ if $X=0$, and admit it to the ordinary verifier if $X=1$. Since $p=q$ an admitted proposal is always accepted, and the output law is
\begin{equation}
\Pr(Y=0)=\tfrac14,\qquad\Pr(Y=1)=\tfrac34.\label{eq:bias}
\end{equation}
The acceptance formula was unchanged; admission selected proposals by their realized value. A mandatory progress floor on the first token does not help, since the same construction applies at the second position. Three settings avoid the issue: greedy decoding; target-only verification with point-mass proposals, where there is no proposal law to bias; and fixed-noise verification (Theorem~\ref{thm:fixednoise}), under which every depth policy yields the same output. The demand frontier above changes only the amount and order of arithmetic used to decide a fixed protocol, so it is exact pointwise.

\subsection{Measurements}\label{sec:decisions-measured}
\pending{geometry, kernel: Gumbel-race candidate counts through the int8 pass on real states at $T=1$ and $T=0.7$; mass-enclosure ratios and $P_?$ for standard acceptance on real MTP and DFlash drafts; the fraction of decisions resolved before any re-scoring}
